%% ============================================================
%% FRONT MATTER -- Applied Energy submission
%% pandoc-friendly LaTeX fragment.
%% GENERATED: the abstract and highlights below are copied verbatim from V6.tex
%% by scripts/sync_frontmatter.py. Edit V6.tex, then re-run that script; the gate
%% checks the two agree, because an earlier divergence shipped a stale abstract in
%% the Word twin.
%% ============================================================

%% ABSTRACT
\section*{Abstract}

Hydrogen for European residential heating is usually assessed at national resolution or imposed as an adoption share, and neither tests whether peak-serving assets recover capital market by market. We reconstruct heat demand on a 1{,}369-region NUTS3 surface across 29 markets, aggregating to the country before each test, over 200 draws per scenario. Four scenarios impose exogenous hydrogen-adoption settings, compared with separately derived winter-peak dispatch support. Hydrogen is tested in three roles: annual levelised cost, cold-snap dispatch across buildings, district heat and power, and peaker capital recovery.

At central 2050 parameters the best heat pump delivers at \euro{}118.5/MWh against \euro{}122.8 for a hydrogen boiler and \euro{}132.7 for gas, beating hydrogen in 22 of 29 markets with firm-generation cost allocated pro rata to heating's demand share, and in 19 to 17 when allocated by the change in peaking capacity. The count spans 13 to 29 across the hydrogen-price paths tested. On the high carbon path and equal-efficiency series, hydrogen undercuts gas at the power peak in a different seven, all salt-cavern. No tested market reaches full peaker-capital recovery over the \euro{}30 to 120/MWh scarcity-margin range; capacity-weighted recovery among the five that build is about 9\%. With the hydrogen-ready turbine at \euro{}600/kW, technology-neutral support favours gas; selecting hydrogen needs an emissions-conditioned premium of about \euro{}172/tCO$_2$ of displaced gas.
Operational residential-heating emissions fall from 644~MtCO$_2$ in 2025 to between 10 and 352 by 2050. Power-sector quantities rest on one constructed weather year and are screening results, not adequacy estimates.

%% HIGHLIGHTS
\section*{Highlights}

\begin{itemize}
\item Heat pumps beat hydrogen on 2050 base-load heat cost in 22 of 29 markets
\item Hydrogen wins the high-carbon power peak in seven cavern markets at equal efficiency
\item No tested market fully recovers peaker capital across the assumed scarcity margins
\item At \euro{}600/kW turbine capital cost, technology-neutral support favours gas
\item Scenario ambition shifts 2050 residential-heating emissions by 342~MtCO$_2$/yr
\end{itemize}

%% KEYWORDS
\section*{Keywords}

Heat pumps; Hydrogen; Building heat decarbonisation; Seasonal hydrogen storage; Missing money

%% NOMENCLATURE
\section*{Nomenclature}

\textit{Latin symbols}
\begin{description}
 \item[$c^{\mathrm{gas}}_{c},\ c^{\mathrm{H_2}}_{c},\ c^{\mathrm{HP}}_{c}$] Short-run marginal cost of delivered heat from the gas boiler, hydrogen boiler and heat pump in country $c$ (\euro{}/MWh-heat)
 \item[$\mathrm{COP}^{\mathrm{cold}}_{c}$] Cold-snap heat-pump coefficient of performance in country $c$
 \item[$\mathrm{CRF}_{c}$] Capital-recovery factor at country discount rate $r_{c}$ over
the asset life (25 yr for the peaking turbine)
 \item[$D_{c}(t)$] National useful heat demand in country $c$, year $t$ (TWh/yr)
 \item[$D_{c}^{0}$] National base-year (2025) useful heat demand in country $c$ (TWh/yr)
 \item[$D^{\mathrm{pk}}$] System peak electricity demand (GW)
 \item[$D^{\mathrm{tot}}$] Total electricity demand (TWh/yr)
 \item[$(\mathrm{Dw}/\mathrm{Pop})_{c,t}$] Dwellings per capita in country $c$, year $t$
 \item[$e_{\mathrm{g}}$] Natural-gas emission factor ($=0.202$ tCO$_2$/MWh)
  \item[$E^{\mathrm{pk}}_{c}$] Annual energy generated by the firm peaking fleet in country $c$ (TWh)
 \item[$f$] Fixed operating-cost rate on peaking-turbine capital ($=2.5\,\%$/yr)
 \item[$\mathrm{FLH}_{c}$] Full-load hours of the hydrogen peaker in country $c$ (h/yr)
 \item[$H_{c}$] Equivalent full-load hours of the peak slice in country $c$ ($\approx$518 h/yr, over a 2{,}000 h window)
 \item[$H_{c}^{2015}$] Hotmaps 2015 benchmark demand for country $c$ (TWh)
 \item[$K^{\mathrm{T}}$] Hydrogen-ready peaking-turbine overnight capital cost ($=$
\euro{}600/kW; conventional open-cycle gas $=$ \euro{}450/kW). Eq.~\eqref{eq:recovery}
turns it into an annual charge by applying $\mathrm{CRF}_{c}+f$; this is the
\euro{}600/kW of the abstract and of Section~\ref{sec:results:missing}
 \item[$K^{\mathrm{pk}}_{c}$] Firm peaking-fleet capacity in country $c$ (GW)
 \item[$N$] Monte Carlo / Saltelli sample size
 \item[$p^{\mathrm{gas}}_{c}$] Retail gas price in country $c$ (\euro{}/MWh)
 \item[$p^{\mathrm{gas,w}}$] Winter (backbone) gas price (\euro{}/MWh)
 \item[$p^{\mathrm{H_2}}_{c}$] Delivered hydrogen price in country $c$ (\euro{}/MWh)
 \item[$p^{\mathrm{H_2,bb}}_{c}$] Backbone hydrogen price in country $c$ (\euro{}/MWh)
 \item[$\mathrm{Pop}_{c,t}$] Projected population of country $c$, year $t$
 \item[$r_{c,s}$] Per-country, per-scenario envelope-renovation rate (\%/yr)
 \item[$r_{c}$] Country weighted-average cost of capital (WACC)
 \item[$R_{c}$] Capital recovery of the hydrogen peaker in country $c$, its scarcity rent
over its annualised capital-and-fixed-operating charge (dimensionless;
Eq.~\eqref{eq:recovery})
 \item[$\mathrm{SCOP}_{c}$] Seasonal coefficient of performance of the heat pump in country $c$
 \item[$S_{T}$] Sobol total-order sensitivity index
 \item[$t,\ t_{0}$] Projection year; base year ($t_{0}=2025$)
 \item[$z_{c},\ z_{\mathrm{EU}},\ z_{c}^{\mathrm{idio}}$] Standard-normal shocks (common and idiosyncratic) for sampling renovation rates
\end{description}

\textit{Greek symbols}
\begin{description}
 \item[$\varphi_{y}$] Carbon price (ETS2) in year $y$ (\euro{}/tCO$_2$)
 \item[$\eta_{\mathrm{g}},\ \eta_{\mathrm{h}}$] Gas-boiler (0.92) and hydrogen-boiler (0.90) efficiencies
 \item[$\eta_{\mathrm{OCGT}},\ \eta_{\mathrm{T}}$] Open-cycle gas turbine and hydrogen turbine electrical efficiencies. The winter-peak merit order that sets the win counts holds both at $0.40$, so that comparison is even-handed. The capital-recovery layer instead lets them evolve, the gas turbine to $0.42$ by 2050 and the hydrogen turbine to $0.44$, or $0.48$ under H2~Push, which assumes an advanced simple-cycle machine; there the hydrogen unit is deliberately given the better turbine
 \item[$\mu_{c,s},\ \sigma_{c,s}$] Mean and standard deviation of the sampled renovation rate $r_{c,s}$
 \item[$\pi^{\mathrm{peak}}_{c}$] Derived scarcity (peak) electricity price in country $c$ (\euro{}/MWh-e)
 \item[$\pi^{\mathrm{s}}$] Scarcity premium (low/central/high $\{30,60,120\}$ \euro{}/MWh; the sensitivity grid adds 90)
 \item[$\rho$] Cross-country shock correlation: renovation rate ($=0.5$),
 fuel prices ($=0.75$); all other sampled parameters are EU-wide scalars, i.e.
 $\rho=1$ implicitly
 \item[$a^{\mathrm{H_2,dist}}_{c}$] Hydrogen last-mile distribution adder in country $c$, charged in the symmetric accounting (median \euro{}13.4/MWh)
 \item[$\sigma_{c}$] Seasonal hydrogen-storage adder, heat side ($\approx$\euro{}50 to 100/MWh-heat)
 \item[$\sigma^{e}_{c}$] Seasonal hydrogen-storage adder, electricity side (\euro{}/MWh-e)
 \item[$\Delta_{c}^{\mathrm{VM}}$] Vintage-matched deviation of the bottom-up backcast from the Hotmaps 2015 benchmark in country $c$
 \item[$\Phi_{\mathrm{cap}},\ \Phi_{\mathrm{en}}$] Firm dispatchable peaking fleet's share of system peak capacity and of generated energy
\end{description}

\textit{Subscripts and superscripts}
\begin{description}
 \item[$c$] Country (region) index
 \item[$s$] Scenario index
 \item[$t,\ t_{0},\ y$] Year; base year; carbon-price year
 \item[$\mathrm{gas},\ \mathrm{H_2},\ \mathrm{HP}$] Gas-boiler, hydrogen-boiler and heat-pump branches
 \item[$\mathrm{pk},\ \mathrm{peak},\ \mathrm{cold}$] Peaker; system-peak/scarcity; cold-snap conditions
\end{description}
